% GENERATED FILE. Edit canonical lessons, then run npm run generate:books.
% canonical-source-hash: fnv1a64:a5e01e6e7a8a40fd
% canonical-lessons: LA-C385-rego, LA-C385-guberno, LA-C385-administro, LA-C385-imperium, LA-C385-princeps

\chapter{Ruling and Governing}
\label{ch:la-ruling-and-governing}
\hlchaptermodality{385}

\begin{chapteropening}
\textbf{By the end of this chapter:} \emph{I can say that someone rules, steers or runs a state, and talk about supreme power and the first man of the state.}

Complete the last lesson of chapter 385: I can say that someone rules, steers or runs a state, and talk about supreme power and the first man of the state.
\end{chapteropening}

\section[regō, regere]{regō, regere — to rule, to guide, to direct}

\label{lesson:LA-C385-rego}



\begin{quote}
Before the new one: say the Latin for canvassing, then the Latin for to pick out.
\end{quote}

\subsection*{You'll want to know: regō, regere}
\textbf{regō, regere} — \textquotedblleft{}to rule, to guide, to direct\textquotedblright{}.

\begin{quote}
Rēx populum iūstē regēbat.
\end{quote}

\begin{quote}
\emph{The king ruled the people justly.}
\end{quote}

\emph{regō, regere, rēxī, rēctum}. \emph{Rēctus}, straight, is its participle; \emph{rēx} is the man who rules. English \emph{regent} and \emph{correct}.

\begin{cousinweb}
\emph{regō} is inherited from a Proto-Indo-European root meaning \textquotedblleft{}to make straight, to direct\textquotedblright{} (EtymDB), the same root as in \emph{rēx} and \emph{rēctus} (Lewis \& Short). English \emph{regent} and \emph{correct} come from it.
\end{cousinweb}

\begin{grammarlens}[title={What you've built}]
One more word for telling, explaining and describing at B1.
\end{grammarlens}

\subsection*{Your turn}
\begin{itemize}
\raggedright
  \item \emph{Say it:} \emph{regō, regere}
  \item \emph{Say it:} \emph{regō, regere}, once more
  \item \emph{Say it:} a sentence of your own with \emph{regō, regere}
  \item \emph{Recall:} say the Latin for canvassing, then the Latin for to pick out, then say \emph{regō, regere} again
\end{itemize}

\begin{tcolorbox}[breakable,colback=teal!4,colframe=teal!35!black,title={Before you move on}]
What does regō, regere mean? (\textquotedblleft{}To rule, to guide, to direct\textquotedblright{}.) Say it once more.
\end{tcolorbox}
\section[gubernō, gubernāre]{gubernō, gubernāre — to steer; to govern}

\label{lesson:LA-C385-guberno}



\begin{quote}
Before the new one: say the Latin for to pick out, then the Latin for to rule.
\end{quote}

\subsection*{You'll want to know: gubernō, gubernāre}
\textbf{gubernō, gubernāre} — \textquotedblleft{}to steer; to govern\textquotedblright{}.

\begin{quote}
Nauta nāvem gubernat.
\end{quote}

\begin{quote}
\emph{The sailor steers the ship.}
\end{quote}

\emph{gubernō, gubernāre, gubernāvī, gubernātum}. \emph{Gubernātor}, a helmsman, is already known. English \emph{govern} comes from it through French.

\begin{cousinweb}
\emph{gubernō} was borrowed from Greek \emph{kybernaō} \textquotedblleft{}to steer a ship\textquotedblright{} (EtymDB; Lewis \& Short). English \emph{govern} comes from it through French.
\end{cousinweb}

\begin{grammarlens}[title={What you've built}]
One more word for telling, explaining and describing at B1.
\end{grammarlens}

\subsection*{Your turn}
\begin{itemize}
\raggedright
  \item \emph{Say it:} \emph{gubernō, gubernāre}
  \item \emph{Say it:} \emph{gubernō, gubernāre}, once more
  \item \emph{Say it:} a sentence of your own with \emph{gubernō, gubernāre}
  \item \emph{Recall:} say the Latin for to pick out, then the Latin for to rule, then say \emph{gubernō, gubernāre} again
\end{itemize}

\begin{tcolorbox}[breakable,colback=teal!4,colframe=teal!35!black,title={Before you move on}]
What does gubernō, gubernāre mean? (\textquotedblleft{}To steer; to govern\textquotedblright{}.) Say it once more.
\end{tcolorbox}
\section[administrō, administrāre]{administrō, administrāre — to manage, to run; to administer}

\label{lesson:LA-C385-administro}



\begin{quote}
Before the new one: say the Latin for to rule, then the Latin for to steer.
\end{quote}

\subsection*{You'll want to know: administrō, administrāre}
\textbf{administrō, administrāre} — \textquotedblleft{}to manage, to run; to administer\textquotedblright{}.

\begin{quote}
Prōvinciam bene administrāvit.
\end{quote}

\begin{quote}
\emph{He governed the province well.}
\end{quote}

\emph{administrō, administrāre, administrāvī, administrātum}, built on \emph{minister}, a servant, a helper. English \emph{administer}.

\begin{cousinweb}
\emph{administrō} is built from \emph{ad} and \emph{ministrō} \textquotedblleft{}to serve\textquotedblright{} (EtymDB), which comes from \emph{minister} \textquotedblleft{}a servant, a helper\textquotedblright{} (Lewis \& Short). English \emph{administer} comes from it.
\end{cousinweb}

\begin{grammarlens}[title={What you've built}]
One more word for telling, explaining and describing at B1.
\end{grammarlens}

\subsection*{Your turn}
\begin{itemize}
\raggedright
  \item \emph{Say it:} \emph{administrō, administrāre}
  \item \emph{Say it:} \emph{administrō, administrāre}, once more
  \item \emph{Say it:} a sentence of your own with \emph{administrō, administrāre}
  \item \emph{Recall:} say the Latin for to rule, then the Latin for to steer, then say \emph{administrō, administrāre} again
\end{itemize}

\begin{tcolorbox}[breakable,colback=teal!4,colframe=teal!35!black,title={Before you move on}]
What does administrō, administrāre mean? (\textquotedblleft{}To manage, to run; to administer\textquotedblright{}.) Say it once more.
\end{tcolorbox}
\section[imperium, imperiī]{imperium, imperiī — command, supreme power; an empire}

\label{lesson:LA-C385-imperium}



\begin{quote}
Before the new one: say the Latin for to steer, then the Latin for to manage.
\end{quote}

\subsection*{You'll want to know: imperium, imperiī}
\textbf{imperium, imperiī} — \textquotedblleft{}command, supreme power; an empire\textquotedblright{}.

\begin{quote}
Cōnsulēs summum imperium habent.
\end{quote}

\begin{quote}
\emph{The consuls hold supreme power.}
\end{quote}

\emph{imperium, imperiī} n., built on \emph{imperō}, to command: the power to command an army and govern, then the land ruled. English \emph{empire} and \emph{imperial}.

\begin{cousinweb}
\emph{imperium} is built on \emph{imperō} \textquotedblleft{}to command\textquotedblright{} (EtymDB; Lewis \& Short). English \emph{empire} comes from it through French, and \emph{imperial} from its adjective.
\end{cousinweb}

\begin{grammarlens}[title={What you've built}]
One more word for telling, explaining and describing at B1.
\end{grammarlens}

\subsection*{Your turn}
\begin{itemize}
\raggedright
  \item \emph{Say it:} \emph{imperium, imperiī}
  \item \emph{Say it:} \emph{imperium, imperiī}, once more
  \item \emph{Say it:} a sentence of your own with \emph{imperium, imperiī}
  \item \emph{Recall:} say the Latin for to steer, then the Latin for to manage, then say \emph{imperium, imperiī} again
\end{itemize}

\begin{tcolorbox}[breakable,colback=teal!4,colframe=teal!35!black,title={Before you move on}]
What does imperium, imperiī mean? (\textquotedblleft{}Command, supreme power; an empire\textquotedblright{}.) Say it once more.
\end{tcolorbox}
\section[prīnceps, prīncipis]{prīnceps, prīncipis — the first, a leader; the emperor}

\label{lesson:LA-C385-princeps}



\begin{quote}
Before the new one: say the Latin for to manage, then the Latin for command.
\end{quote}

\subsection*{You'll want to know: prīnceps, prīncipis}
\textbf{prīnceps, prīncipis} — \textquotedblleft{}the first, a leader; the emperor\textquotedblright{}.

\begin{quote}
Prīnceps cīvitātis erat.
\end{quote}

\begin{quote}
\emph{He was the leading man of the state.}
\end{quote}

\emph{prīnceps, prīncipis} m.: \emph{prīmus} + \emph{capiō}, the one who takes first place. As an adjective it means first. English \emph{prince} and \emph{principal}.

\begin{cousinweb}
\emph{prīnceps} is built from \emph{prīmus} \textquotedblleft{}first\textquotedblright{} and \emph{capiō} \textquotedblleft{}to take\textquotedblright{} (EtymDB; Lewis \& Short): the one who takes first place. English \emph{prince} and \emph{principal} come from it.
\end{cousinweb}

\begin{grammarlens}[title={What you've built}]
One more word for telling, explaining and describing at B1.
\end{grammarlens}

\subsection*{Your turn}
\begin{itemize}
\raggedright
  \item \emph{Say it:} \emph{prīnceps, prīncipis}
  \item \emph{Say it:} \emph{prīnceps, prīncipis}, once more
  \item \emph{Say it:} a sentence of your own with \emph{prīnceps, prīncipis}
  \item \emph{Recall:} say the Latin for to manage, then the Latin for command, then say \emph{prīnceps, prīncipis} again
\end{itemize}

\begin{tcolorbox}[breakable,colback=teal!4,colframe=teal!35!black,title={Before you move on}]
What does prīnceps, prīncipis mean? (\textquotedblleft{}The first, a leader; the emperor\textquotedblright{}.) Say it once more.
\end{tcolorbox}
